\begin{abstract}
Existing LLM benchmarks for EDA target single-function code generation (RTLLM, VerilogEval) or open-source flow repair (FlowFixer-EDA), but none evaluate the end-to-end \emph{orchestration} of multi-step design flows with live commercial EDA tools.
% 现有的面向 EDA 的 LLM 基准测试只关注单函数代码生成（RTLLM、VerilogEval）或开源流程修复（FlowFixer-EDA），尚无一个基准真正评估在真实商业 EDA 工具上对多步设计流进行端到端"编排"的能力。
We present \bench, a benchmark of \textbf{\Ntasks{} tasks} (Full set) across \textbf{13 EDA domains}---from schematic/netlist to SPICE simulation, physical verification, and digital synthesis---each derived from \emph{real coding-agent traces} of IC design engineers in a production environment, with a \textbf{\LiteSet{}-task Lite set} released for efficient external evaluation.
% 我们提出 \bench，包含 \textbf{\Ntasks{} 个任务}（Full 集），覆盖 \textbf{13 个 EDA 领域}——从原理图/网表到 SPICE 仿真、物理验证和数字综合——每一个任务都来源于真实生产环境中 IC 设计工程师使用编程智能体的交互轨迹；同时发布 \textbf{\LiteSet{} 个任务的 Lite 集}用于高效的外部评测。
Each task requires the agent to produce a deliverable package (\texttt{submit.zip}) by orchestrating remote EDA tools through an MCP bridge, scored against $\geq 8$ executable assertions with a \texttt{PreToolUse} hook that prevents answer leakage.
% 每个任务要求智能体通过 MCP 网桥调用远程 EDA 工具，最终产出一个交付件包（\texttt{submit.zip}）；判分基于不少于 8 条可执行断言，并通过 \texttt{PreToolUse} 钩子防止答案泄露。
On a frozen \Neval{}-run evaluation snapshot, a single frontier backbone (DeepSeek-V4-Pro) reaches only \textbf{10.0\% Pass@1} (62/620 decided runs)---versus 76.5\% on single-script EDA benchmarks---showing that multi-step flow orchestration is qualitatively harder than isolated code generation.
% 在冻结的 \Neval{}-run 评测快照上，单一前沿骨干（DeepSeek-V4-Pro）仅取得 \textbf{10.0\% Pass@1}（62/620 已判定 run）——而单脚本 EDA 基准上为 76.5\%——表明多步流程编排比孤立的代码生成在本质上更为困难。
Failure analysis attributes the gap to \emph{tool-chain grounding depth}: the highest-passing domain (pnr\_signoff, 20.8\%) exceeds the lowest (physical\_verify and layout\_aether, both 2.9\%) by a factor of $7.1\times$ (a descriptive point estimate; the two confidence intervals overlap), and the lowest-passing domains are precisely those demanding the deepest live-tool interaction (DRC/LVS runs, GUI-driven layout editing).
% 失败分析把差距归因于"工具链接地深度"：通过率最高的领域（pnr\_signoff，20.8\%）与最低的领域（physical\_verify 与 layout\_aether，均为 2.9\%）相差 7.1 倍（描述性点估计；两个置信区间重叠），而通过率最低的领域恰好是需要最深实时工具交互（DRC/LVS 运行、GUI 驱动的版图编辑）的领域。
\end{abstract}
